\section{終わりに}
今回の内容はどうだっただろうか？テスト範囲のみの説明を書きつつ,それに付随する(発展)内容を記載していくと,どこまで書けばいいのかわからず,
結局作成時の自分の知っている内容の全てに近い量が掲載されている.特にSection.22は「応用線形代数学」に入りつつあり,個人的には実践的な数学の方が面白いと感じるので,2a,2bが応用に近い分,かえって構成が難しかった.
下記には参考文献があり,\cite{physics}に関しては大学院向けのため,購入するには早とちり過ぎた.しかし,面白そうな内容も多くペラペラとたまに読んでいる.
前期と異なり,線形代数学は紙面上の数学に限らないより実践的な背景を感じることができたのではないかと考えている.2年になっても線形代数学は使うので,積極的に復習や勉強をしてほしいと考えている.\\
\hspace{1em}
最後に,このテキストの作成は私が大した知識も無い中で,純粋に面白いと思った内容をみんなに共有したいという気持ちから作成が始まった.結果として,読者の成績に影響を与えることができて素直に嬉しく感じる.
そして,読んでくださったことに感謝している.私の線形代数学の知識など,面白いと思ったものをかいつまんでいるに過ぎず,大したことはない.そのため,これからも線形代数学は興味のままに学習しようと考えている.実は,線形代数学の知識量は読者とそう変わらない.
一緒に学習していければと思っている.再度となるが,今までありがとうございました.
\section{参考文献}
\begin{thebibliography}{9}
  \bibitem{strang}
  Gilbert Strang (松崎公紀(他)訳)『ストラング：線形代数イントロダクション(原著 第4版)』近代科学社, 2022.
  \bibitem{physics}
  Gilbert Strang (日本応用数理学会(他)訳)『ストラング：計算理工学』近代科学社, 2019.
  \bibitem{calclus}
  James Stewart (秋山仁・伊藤雄二 訳)『スチュワート微分積分学Ⅲ 多変数関数の微積分』
  \bibitem{explain}
  関口陽太 他「2024年度線形代数学(1a), (1b), (2a), (2b)解答解説」,4bitcom過去問制作委員会, 2024-2025.
  \bibitem{lenar_pre}
  関口陽太『前期線形代数学Ver.2.10』, 2025.
  \bibitem{lenar_post}
  関口陽太『後期線形代数学Ver.1.00』, 2025.
  \bibitem{Trefethen-Bau}
  Lloyd N. Trefethen and David Bau III『Numerical Linear Algebra』SIAM, 2022.
\end{thebibliography}
\newpage